\chapter{Limits and Continuity}

本章讨论 sequence limit 和 function limit。两者都在描述“靠近”：sequence limit
说的是当 \(n\) 足够大时，\(x_n\) 靠近 \(L\)；function limit 说的是当
\(x\) 足够靠近 \(a\) 但不等于 \(a\) 时，\(f(x)\) 靠近 \(L\)。本章的重点不是
记住某个直观图像，而是把“足够靠近”翻译成可以检查的
\(\varepsilon\)-\(N\) 或 \(\varepsilon\)-\(\delta\) inequality proof。

\section{From Sequence Limits to Function Limits}

对 sequence \((x_n)\)，
\[
  \lim_{n\to\infty}x_n=L
\]
的意思是
\[
  \forall \varepsilon>0,\ \exists N\in\N
  \text{ such that } n\geq N \Longrightarrow |x_n-L|<\varepsilon.
\]
这里 \(\varepsilon\) 是允许的 output error，而 \(N\) 是从哪一项开始进入这个
error bound。证明 sequence limit 时，通常先从
\(|x_n-L|<\varepsilon\) 倒推需要 \(n\) 多大，再把这个要求写成一个合法的
\(N\)。

Function limit 的输入不是离散的 \(n=0,1,2,\ldots\)，而是实数 \(x\) 从
一段 interval 内靠近某一点 \(a\)。因此 \(N\) 被替换成 \(\delta\)：\(\delta\)
控制 input distance，\(\varepsilon\) 控制 output distance。
这两个 definitions 的 quantifier order 完全平行：证明者先听到任意给定的
error tolerance，再回应一个足够严格的 input tolerance。不能先固定一个
\(\delta\) 或 \(N\) 再要求它适合所有 \(\varepsilon\)，因为越小的
\(\varepsilon\) 需要越精细的控制。

\begin{definition}[function limit]
设 \(I\) 是包含 \(a\) 的 open interval，且
\[
  f:I\setminus\{a\}\to\R.
\]
我们说
\[
  \lim_{x\to a}f(x)=L
\]
当且仅当
\[
  \forall \varepsilon>0,\ \exists \delta>0
  \text{ such that }
  0<|x-a|<\delta \Longrightarrow |f(x)-L|<\varepsilon.
\]
\end{definition}

条件 \(0<|x-a|\) 的作用是排除 \(x=a\)。Function limit 关心的是
\(x\) approaching \(a\) 时附近的 function values，而不是 \(f(a)\) 本身。
因此一个 function 可以在 \(a\) 没有定义，但仍然在 \(a\) 有 limit。

\section{First Epsilon-Delta Proofs}

最基本的 \(\varepsilon\)-\(\delta\) proof 有固定格式。给定任意
\(\varepsilon>0\)，我们要选择一个 \(\delta>0\)。选择 \(\delta\) 之前，
先把 \(|f(x)-L|\) 改写成含有 \(|x-a|\) 的表达式。改写完成后，选择
\(\delta\) 使得
\[
  0<|x-a|<\delta
\]
能够推出
\[
  |f(x)-L|<\varepsilon.
\]

\begin{reviewbox}{Epsilon-delta proof template}
写 \(\varepsilon\)-\(\delta\) proof 时，可以按四步组织。第一步，写清楚
``Let \(\varepsilon>0\) be given''，并声明目标是找到 \(\delta>0\)。第二步，
在 scratch work 中从 \(|f(x)-L|\) 出发，把它化成
``constant \(\times |x-a|\)'' 或者
``bounded factor \(\times |x-a|\)''。第三步，若表达式中有 denominator,
square root, 或其它可能变化的 factor，先用一个小的固定条件如
\(\delta\leq1\) 或 \(\delta\leq1/2\) 把这些 factor bounded。第四步，把所有
条件合并成
\[
  \delta=\min\{\text{固定控制},\varepsilon\text{-控制}\}.
\]
正式 proof 中不需要展示所有倒推过程，但必须展示 chosen \(\delta\) 如何让
每个 inequality 成立。
\end{reviewbox}

例如
\[
  \lim_{x\to3}(2x+1)=7.
\]
因为
\[
  |(2x+1)-7|=|2x-6|=2|x-3|,
\]
所以只要让 \(|x-3|<\varepsilon/2\)，就能让 output error 小于
\(\varepsilon\)。于是可以取 \(\delta=\varepsilon/2\)。

\begin{exerciseblock}[Lecture Notes Exercise 57]
已知
\[
  \lim_{x\to -2}(4x+3)=-5.
\]
当 \(\varepsilon=0.1\) 时，找出一个可用的 \(\delta\)。
\end{exerciseblock}

\begin{solution}
这里
\[
  f(x)=4x+3,\qquad a=-2,\qquad L=-5.
\]
我们计算 output error：
\[
  |f(x)-L|
  =
  |(4x+3)-(-5)|
  =
  |4x+8|
  =
  4|x+2|.
\]
现在要求
\[
  |f(x)-L|<0.1.
\]
由上式可知，只要
\[
  4|x+2|<0.1,
\]
也就是
\[
  |x+2|<0.025,
\]
就满足要求。因此取
\[
  \delta=0.025.
\]
若 \(0<|x-(-2)|<\delta\)，即 \(0<|x+2|<0.025\)，则
\[
  |(4x+3)-(-5)|
  =
  4|x+2|
  <
  4(0.025)
  =
  0.1.
\]
所以 \(\delta=0.025\) 是一个符合 \(\varepsilon=0.1\) 的选择。
\end{solution}

\begin{explanation}
这道题不是要求重新证明所有 \(\varepsilon>0\) 的情况，而是在已知 limit
存在的背景下，为指定的 tolerance 找一个 input tolerance。线性函数的误差
正好是 \(4|x+2|\)，所以 \(\delta\) 只需要把这个四倍关系抵消掉。任何
\(0<\delta\leq 0.025\) 也都可以使用；选择最大这个数只是最直接。
遇到 affine function \(mx+b\) 的同类问题，先计算
\[
  |(mx+b)-(ma+b)|=|m|\,|x-a|.
\]
若题目给定一个具体 \(\varepsilon_0\)，就取
\(\delta\leq \varepsilon_0/|m|\)。若题目要求证明所有
\(\varepsilon>0\)，就把同一个式子写成
\(\delta=\varepsilon/|m|\)；当 \(m=0\) 时 function 是 constant，任意正
\(\delta\) 都可用。
\end{explanation}

\begin{example}[constant function]
若 \(f(x)=c\)，则对任意 \(a\)，
\[
  \lim_{x\to a}c=c.
\]
给定 \(\varepsilon>0\)，可取 \(\delta=1\)。只要
\(0<|x-a|<1\)，就有
\[
  |f(x)-c|=|c-c|=0<\varepsilon.
\]
\end{example}

\begin{example}[a function with a hole]
考虑
\[
  f:(0,\infty)\setminus\{2\}\to\R,\qquad
  f(x)=\frac{x^2-4}{x-2}.
\]
当 \(x\neq2\) 时，
\[
  f(x)=\frac{(x-2)(x+2)}{x-2}=x+2.
\]
因此
\[
  \left|\frac{x^2-4}{x-2}-4\right|
  =
  |(x+2)-4|
  =
  |x-2|.
\]
给定 \(\varepsilon>0\)，取 \(\delta=\varepsilon\)。若
\(0<|x-2|<\delta\)，则
\[
  \left|\frac{x^2-4}{x-2}-4\right|
  =
  |x-2|
  <
  \varepsilon.
\]
所以
\[
  \lim_{x\to2}\frac{x^2-4}{x-2}=4.
\]
这个例子说明 limit 不需要 \(f(2)\) 有定义。
\end{example}

\begin{exerciseblock}[Lecture Notes Exercise 58]
证明
\[
  \lim_{x\to4}\sqrt{x}=2.
\]
\end{exerciseblock}

\begin{solution}
函数 \(\sqrt{x}\) 的 natural domain 是 \(x\geq0\)。在 \(x\) 靠近 \(4\)
时，相关的 \(x\) 都可以限制在正数范围内。给定 \(\varepsilon>0\)，我们先估计
\[
  |\sqrt{x}-2|.
\]
对 \(x\geq0\)，有
\[
  |\sqrt{x}-2|
  =
  \left|\frac{(\sqrt{x}-2)(\sqrt{x}+2)}{\sqrt{x}+2}\right|
  =
  \frac{|x-4|}{\sqrt{x}+2}.
\]
因为 \(\sqrt{x}+2\geq2\)，所以
\[
  |\sqrt{x}-2|
  \leq
  \frac{|x-4|}{2}.
\]
现在取
\[
  \delta=2\varepsilon.
\]
若 \(0<|x-4|<\delta\) 且 \(x\geq0\)，则
\[
  |\sqrt{x}-2|
  =
  \frac{|x-4|}{\sqrt{x}+2}
  \leq
  \frac{|x-4|}{2}
  <
  \frac{\delta}{2}
  =
  \varepsilon.
\]
因此
\[
  \lim_{x\to4}\sqrt{x}=2.
\]
\end{solution}

\begin{explanation}
困难点在于 \(\sqrt{x}-2\) 不直接含有 \(|x-4|\)。乘以 conjugate
\(\sqrt{x}+2\) 后，分子变成 \(x-4\)，这正好连接到 input distance。
分母 \(\sqrt{x}+2\) 不会接近 \(0\)；它至少是 \(2\)，所以可以用
\(|x-4|/2\) 控制整个 error。这里不需要让 \(\delta\leq1\)，因为分母的下界
已经由 \(x\geq0\) 给出。
解决类似 radical limit 时，先寻找能把 radical difference 变成 polynomial
difference 的 conjugate。一般地，
\[
  |\sqrt{x}-\sqrt{a}|=\frac{|x-a|}{\sqrt{x}+\sqrt{a}}
\]
在 \(a>0\) 时会有分母下界；在 \(a=0\) 时则不能用同一套 estimate，需要从
\(|\sqrt{x}|<\varepsilon\) 倒推出 \(|x|<\varepsilon^2\)。因此方法不是记住
某个 \(\delta\)，而是先判断 denominator 是否能在 limiting point 附近远离
\(0\)。
\end{explanation}

\section{When Limits Do Not Exist}

证明某个给定的 \(L\) 不是 limit，需要否定 \(\varepsilon\)-\(\delta\)
definition：
\[
  \lim_{x\to a}f(x)\neq L
\]
等价于存在一个 \(\varepsilon>0\)，使得对每个 \(\delta>0\)，都能找到
某个 \(x\) 满足
\[
  0<|x-a|<\delta
  \quad\text{and}\quad
  |f(x)-L|\geq\varepsilon.
\]
证明 limit does not exist 更强：我们要说明没有任何 \(L\in\R\) 能通过这个
definition。常用方法是找两列靠近同一点的 inputs，但对应的 function values
靠近不同的 numbers。

\begin{reviewbox}{How to prove nonexistence of a function limit}
有两种常见路线。第一种是直接使用 negation of the definition：选定一个
\(\varepsilon_0>0\)，证明每个 \(\delta>0\) 都允许某个
\(0<|x-a|<\delta\) 使 \(|f(x)-L|\geq\varepsilon_0\)。这通常用于否定某个
指定的 candidate \(L\)。第二种是 sequential test：构造
\(x_n\to a\) and \(y_n\to a\)，但 \(f(x_n)\to A\)、\(f(y_n)\to B\)，其中
\(A\neq B\)。这同时排除所有 candidate limits，因为同一个 function limit
会迫使所有 approaching sequences 的 images 收敛到同一个 number。
\end{reviewbox}

\begin{example}
定义
\[
  f(x)=\frac{|x|}{x}
  =
  \begin{cases}
    1, & x>0,\\
    -1, & x<0.
  \end{cases}
\]
当 \(x\to0^+\) 时，function value 恒为 \(1\)；当 \(x\to0^-\) 时，
function value 恒为 \(-1\)。任何一个 real number \(L\) 都不能同时让
\(1\) 和 \(-1\) 在任意小的 tolerance 内都靠近它。因此
\[
  \lim_{x\to0}\frac{|x|}{x}
\]
does not exist。
\end{example}

\begin{exerciseblock}[Lecture Notes Exercise 59]
设 \(f:\R\setminus\{0\}\to\R\) 由
\[
  f(x)=\sin(1/x)
\]
给出。证明 \(\lim_{x\to0}f(x)\) does not exist。
\end{exerciseblock}

\begin{solution}
我们证明不存在 \(L\in\R\) 使得 \(\lim_{x\to0}\sin(1/x)=L\)。

对每个 \(n\in\N\) 且 \(n\geq1\)，定义
\[
  x_n=\frac{1}{2\pi n+\pi/2},
  \qquad
  y_n=\frac{1}{2\pi n+3\pi/2}.
\]
由于分母趋向 \(+\infty\)，有
\[
  x_n\to0,\qquad y_n\to0.
\]
并且每一项都为正数，所以它们都属于 \(\R\setminus\{0\}\)。代入函数得到
\[
  f(x_n)=\sin(2\pi n+\pi/2)=1,
  \qquad
  f(y_n)=\sin(2\pi n+3\pi/2)=-1.
\]

现在假设某个 \(L\in\R\) 是 limit。取 \(\varepsilon=1\)。若
\(\lim_{x\to0}f(x)=L\)，则存在 \(\delta>0\)，使得
\[
  0<|x|<\delta \Longrightarrow |\sin(1/x)-L|<1.
\]
因为 \(x_n\to0\) 且 \(y_n\to0\)，可以取足够大的 \(n\)，使得
\[
  0<x_n<\delta,\qquad 0<y_n<\delta.
\]
于是应当同时有
\[
  |1-L|<1,\qquad |-1-L|<1.
\]
但 triangle inequality 给出
\[
  2=|1-(-1)|
  =
  |(1-L)+(L+1)|
  \leq
  |1-L|+|L+1|
  <
  1+1
  =
  2,
\]
这推出 \(2<2\)，矛盾。因此不存在这样的 \(L\)，所以
\[
  \lim_{x\to0}\sin(1/x)
\]
does not exist。
\end{solution}

\begin{explanation}
这道题的难点是 limit nonexistence 不能只说函数“震荡”。证明必须制造出
definition 无法容纳的冲突。两列 inputs 都趋向 \(0\)，但 outputs 分别固定为
\(1\) 和 \(-1\)。若同一个 \(L\) 是 limit，那么足够靠近 \(0\) 的这两类
outputs 都必须落在半径 \(1\) 的 interval \((L-1,L+1)\) 中；这个 interval
长度为 \(2\)，不能同时严格包含距离为 \(2\) 的两个端点 \(1\) 和 \(-1\)。
做类似 \(f(x)=\sin(1/x)\), \(\cos(1/x)\), 或 sign-changing oscillation
的题时，先选让 inner angle 命中特殊值的 inputs。这里我们不是随机取
sequence，而是解方程 \(1/x=2\pi n+\pi/2\) 和
\(1/x=2\pi n+3\pi/2\)。这种做法同时保证 \(x\to0\) 和 outputs 固定在两
个互不相容的值。
\end{explanation}

\section{Limit Laws}

Limit laws 允许我们把已经证明过的 limits 组合起来。设
\[
  \lim_{x\to a}f(x)=L,\qquad
  \lim_{x\to a}g(x)=M.
\]
则
\[
  \lim_{x\to a}(f(x)+g(x))=L+M,
\]
\[
  \lim_{x\to a}(f(x)g(x))=LM,
\]
并且当 \(M\neq0\) 时，
\[
  \lim_{x\to a}\frac{f(x)}{g(x)}=\frac{L}{M}.
\]

Sum law 的 proof 使用 triangle inequality：
\[
  |(f(x)+g(x))-(L+M)|
  \leq
  |f(x)-L|+|g(x)-M|.
\]
给定 \(\varepsilon>0\)，分别用 \(\varepsilon/2\) 控制两项，再取两个
\(\delta\) 的 minimum。

Product law 需要一个额外步骤：先由 \(f(x)\to L\) 得到 \(f(x)\) 在 \(a\)
附近 bounded。具体地，先让 \(|f(x)-L|<1\)，于是
\[
  |f(x)|\leq |L|+1.
\]
然后写
\[
  |f(x)g(x)-LM|
  \leq
  |f(x)|\,|g(x)-M|+|M|\,|f(x)-L|.
\]
第一项由 \(g(x)\to M\) 控制，第二项由 \(f(x)\to L\) 控制。

Quotient law 的核心是 denominator。若 \(M\neq0\)，先让
\[
  |g(x)-M|<\frac{|M|}{2}.
\]
这推出
\[
  |g(x)|>\frac{|M|}{2},
\]
所以 \(g(x)\) 在 \(a\) 附近不会等于 \(0\)，且 \(1/g(x)\) 的大小可控。
之后把 quotient 写成 product 来处理。

\begin{reviewbox}{Using limit laws without hiding the proof}
Limit laws 是 theorem，不是 algebra shortcut。使用它们之前要确认每个
component limit 都已经存在，并且 quotient 的 limiting denominator
不是 \(0\)。如果题目要求 ``不用 limit laws''，就回到
\(\varepsilon\)-\(\delta\) proof：先把目标 expression 与 candidate limit
相减，通分或展开，直到出现一个会趋向 \(0\) 的 factor，再控制剩余 factors。
如果题目问某个 law 是否能反向使用，优先寻找 cancellation example，因为
两个不存在的 limits 可能在 sum 或 product 中互相抵消。
\end{reviewbox}

\begin{exerciseblock}[Lecture Notes Exercise 60]
找一对 functions，使得
\[
  \lim_{x\to a}f(x)
  \quad\text{and}\quad
  \lim_{x\to a}g(x)
\]
都不存在，但
\[
  \lim_{x\to a}(f(x)+g(x))
\]
存在。
\end{exerciseblock}

\begin{solution}
取 \(a=0\)，定义在 \(\R\setminus\{0\}\) 上的 functions
\[
  f(x)=\frac{|x|}{x},
  \qquad
  g(x)=-\frac{|x|}{x}.
\]
也就是
\[
  f(x)=
  \begin{cases}
    1, & x>0,\\
    -1, & x<0,
  \end{cases}
  \qquad
  g(x)=
  \begin{cases}
    -1, & x>0,\\
    1, & x<0.
  \end{cases}
\]

先看 \(f\)。令
\[
  u_n=\frac1n,\qquad v_n=-\frac1n.
\]
则 \(u_n\to0\) 且 \(v_n\to0\)，但
\[
  f(u_n)=1,\qquad f(v_n)=-1.
\]
若某个 \(L\in\R\) 是 \(\lim_{x\to0}f(x)\)，那么取 \(\varepsilon=1\) 后，
足够大的 \(n\) 会同时满足
\[
  |f(u_n)-L|<1,\qquad |f(v_n)-L|<1.
\]
这会给出
\[
  2=|f(u_n)-f(v_n)|
  \leq
  |f(u_n)-L|+|L-f(v_n)|
  <
  2,
\]
产生矛盾。因此 \(\lim_{x\to0}f(x)\) does not exist。

对 \(g\)，同两列 inputs 给出
\[
  g(u_n)=-1,\qquad g(v_n)=1.
\]
如果 \(\lim_{x\to0}g(x)=M\)，取 \(\varepsilon=1\) 后同样会推出
\[
  2=|g(u_n)-g(v_n)|<2,
\]
这也是矛盾。因此 \(\lim_{x\to0}g(x)\) does not exist。

但是对每个 \(x\neq0\)，
\[
  f(x)+g(x)=\frac{|x|}{x}-\frac{|x|}{x}=0.
\]
因此给定任意 \(\varepsilon>0\)，取 \(\delta=1\)。若
\(0<|x|<\delta\)，则
\[
  |(f(x)+g(x))-0|=|0|=0<\varepsilon.
\]
所以
\[
  \lim_{x\to0}(f(x)+g(x))=0.
\]
\end{solution}

\begin{explanation}
这个例子说明 limit laws 不能反向使用。Sum law 的方向是：若两个 summands
各自有 limit，则 sum 有 limit。反过来，sum 有 limit 并不强迫每个 summand
都有 limit。这里 \(g=-f\) 把 \(f\) 的所有震荡完全抵消，所以 sum 变成
constant zero。
解决类似题目时，先决定想让 sum 变成什么简单 function。若目标是 constant
zero，可以直接取 \(g=-f\)。剩下的任务是挑一个没有 two-sided limit 的
\(f\)，例如 sign function \(|x|/x\) 或 oscillatory function
\(\sin(1/x)\)。然后必须分别证明 \(f\) 和 \(g\) 没有 limit，不能只展示
sum 的计算。
\end{explanation}

\begin{exerciseblock}[Lecture Notes Exercise 61]
不用 limit laws，直接用 \(\varepsilon\)-\(\delta\) proof 计算并证明
\[
  \lim_{x\to0}\frac{x^2-4}{x-1}=4.
\]
\end{exerciseblock}

\begin{solution}
令
\[
  h(x)=\frac{x^2-4}{x-1}.
\]
我们需要证明 \(\lim_{x\to0}h(x)=4\)。给定 \(\varepsilon>0\)，先计算
error：
\[
  \left|\frac{x^2-4}{x-1}-4\right|
  =
  \left|\frac{x^2-4-4x+4}{x-1}\right|
  =
  \left|\frac{x(x-4)}{x-1}\right|.
\]
这里要同时控制 numerator 和 denominator。若 \(|x|<1/2\)，则
\[
  |x-1|\geq 1-|x|>1-\frac12=\frac12,
\]
并且
\[
  |x-4|\leq |x|+4<\frac12+4=\frac92.
\]
因此在 \(|x|<1/2\) 时，
\[
  \left|\frac{x(x-4)}{x-1}\right|
  =
  \frac{|x|\,|x-4|}{|x-1|}
  <
  \frac{|x|(9/2)}{1/2}
  =
  9|x|.
\]
现在取
\[
  \delta=\min\left\{\frac12,\frac{\varepsilon}{9}\right\}.
\]
若 \(0<|x-0|<\delta\)，则 \(|x|<1/2\)，所以上面的 estimate 可用，并且
\[
  \left|\frac{x^2-4}{x-1}-4\right|
  <
  9|x|
  <
  9\delta
  \leq
  \varepsilon.
\]
由 \(\varepsilon\)-\(\delta\) definition，
\[
  \lim_{x\to0}\frac{x^2-4}{x-1}=4.
\]
\end{solution}

\begin{explanation}
Rational function 的 direct proof 通常要先处理 denominator。这里
\(x\to0\)，而 denominator \(x-1\) 靠近 \(-1\)，所以它不会接近 \(0\)。
选择 \(\delta\leq1/2\) 的目的就是把 \(|x-1|\) 固定在大于 \(1/2\) 的范围。
之后 numerator 中剩下一个 \(|x|\)，它会随 \(x\to0\) 变小；另一个因子
\(|x-4|\) 在小 interval 内被 \(9/2\) 控制。于是整个 error 被 \(9|x|\)
控制，再用 \(\delta\leq\varepsilon/9\) 完成证明。
类似 rational function 的 proof 可以先写成一个 checklist：candidate
limit 由代入分子分母得到；error 通分后找出 \(|x-a|\)；对 denominator
设定远离 \(0\) 的固定条件；对其它 bounded factors 设定粗略上界；最后由
\(\varepsilon\) 决定 \(\delta\) 的可变部分。固定条件不必最优，只要能让
后续 inequality 完整成立。
\end{explanation}

\section{Sequential Characterization and Continuity}

\begin{theorem}[sequential characterization of function limits]
Function limit 也可以用 sequences 来刻画：
\[
  \lim_{x\to a}f(x)=L
\]
当且仅当对每个满足
\[
  x_n\in I\setminus\{a\},\qquad x_n\to a
\]
的 sequence，都有
\[
  f(x_n)\to L.
\]
\end{theorem}

\begin{proof}
先证明 forward direction。假设 \(\lim_{x\to a}f(x)=L\)，并取任意
sequence \(x_n\in I\setminus\{a\}\) with \(x_n\to a\)。给定
\(\varepsilon>0\)，由 function limit 可取 \(\delta>0\)，使得
\[
  0<|x-a|<\delta \Longrightarrow |f(x)-L|<\varepsilon.
\]
又因为 \(x_n\to a\)，存在 \(N\) 使得 \(n\geq N\) 时
\(|x_n-a|<\delta\)。每个 \(x_n\neq a\)，所以此时
\[
  |f(x_n)-L|<\varepsilon.
\]
这正是 \(f(x_n)\to L\)。

再证明 reverse direction，使用 contrapositive。若
\(\lim_{x\to a}f(x)\neq L\)，则存在 \(\varepsilon_0>0\)，使得对每个
\(\delta>0\) 都能找到 \(x\in I\setminus\{a\}\) with
\[
  0<|x-a|<\delta
  \quad\text{and}\quad
  |f(x)-L|\geq\varepsilon_0.
\]
对每个 \(n\geq1\)，令 \(\delta=1/n\)，选择这样的 \(x_n\)。于是
\(0<|x_n-a|<1/n\)，所以 \(x_n\to a\)，但
\(|f(x_n)-L|\geq\varepsilon_0\) 对所有 \(n\) 成立，因此
\(f(x_n)\) 不收敛到 \(L\)。这与 sequential condition 矛盾。故 reverse
direction 成立。
\end{proof}

这个 theorem 的用途有两个。若要证明 limit 存在，可以证明所有靠近
\(a\) 的 sequences 都被送到同一个 \(L\)。若要证明 limit 不存在，可以找两列
都趋向 \(a\) 的 inputs，使得对应 outputs 有不同 limits。第二种用法在
oscillation、jump discontinuity、piecewise definition 的题目中最常见。
选择 sequence 时应当让 inputs 自动满足 domain restriction，并让 outputs
落在容易比较的 fixed values 或 known limits。

Continuity 把 limit 和 function value 接起来。若 \(f:I\to\R\) 且
\(a\in I\)，则 \(f\) is continuous at \(a\) 的意思是
\[
  \lim_{x\to a}f(x)=f(a).
\]
用 \(\varepsilon\)-\(\delta\) 写，就是
\[
  \forall\varepsilon>0,\ \exists\delta>0
  \text{ such that }
  |x-a|<\delta \Longrightarrow |f(x)-f(a)|<\varepsilon.
\]
这里不再排除 \(x=a\)，因为 continuity 必须同时检查 \(f(a)\) 本身。
若 function 在 \(a\) 没有定义，它不能 continuous at \(a\)，但它仍可能有
limit。

Continuity proof 与 function limit proof 的差别只有 target value。Limit
proof 的 target \(L\) 可以先猜；continuity proof 的 target 被固定为
\(f(a)\)。因此证明 continuity at \(a\) 时，先计算
\(|f(x)-f(a)|\)，再用与 limit proof 相同的 inequality method 控制它。
证明 discontinuity 时，可以证明 limit 不存在，也可以证明 limit 存在但不等
于 \(f(a)\)，还可以证明 \(f(a)\) 没有定义。

\section{Homework 9}

\begin{exerciseblock}[Homework 9 Problem 1]
证明 discrete limit
\[
  \lim_{n\to\infty}\frac{n^2}{4^n}=0.
\]
\end{exerciseblock}

\begin{solution}
我们使用 \(\varepsilon\)-\(N\) definition。给定 \(\varepsilon>0\)，需要找到
\(N\in\N\)，使得 \(n\geq N\) 时
\[
  \left|\frac{n^2}{4^n}-0\right|<\varepsilon.
\]

先证明一个 estimate：对所有 \(n\geq4\)，
\[
  n^2\leq2^n.
\]
当 \(n=4\) 时，\(4^2=16=2^4\)。若某个 \(n\geq4\) 满足
\(n^2\leq2^n\)，则
\[
  (n+1)^2
  =
  n^2\left(1+\frac1n\right)^2
  \leq
  n^2\left(\frac54\right)^2
  =
  \frac{25}{16}n^2
  <
  2n^2
  \leq
  2^{n+1}.
\]
所以 by induction，\(n^2\leq2^n\) 对所有 \(n\geq4\) 成立。

因此当 \(n\geq4\) 时，
\[
  0\leq\frac{n^2}{4^n}
  \leq
  \frac{2^n}{4^n}
  =
  \frac{1}{2^n}.
\]
由 Archimedean property，选择 \(N\in\N\)，使得
\[
  N\geq4
  \quad\text{and}\quad
  N>\frac1\varepsilon.
\]
当 \(n\geq N\) 时，\(2^n\geq n\geq N>1/\varepsilon\)，所以
\[
  \frac1{2^n}<\varepsilon.
\]
于是
\[
  \left|\frac{n^2}{4^n}-0\right|
  =
  \frac{n^2}{4^n}
  \leq
  \frac1{2^n}
  <
  \varepsilon.
\]
由 \(\varepsilon\)-\(N\) definition，
\[
  \lim_{n\to\infty}\frac{n^2}{4^n}=0.
\]
\end{solution}

\begin{explanation}
这题的核心是 exponential growth eventually dominates polynomial growth。
为了把这个直观变成 proof，我们不用直接处理 \(n^2/4^n\)，而是先证明
\(n^2\leq2^n\)。这样原式被压到 \(1/2^n\) 以下。最后选择 \(N>1/\varepsilon\)
是为了保证 \(1/2^n<\varepsilon\)。每一步都服务于
\(\varepsilon\)-\(N\) definition 中的目标 inequality。
以后遇到 \(\frac{\text{polynomial}}{\text{exponential}}\) 的 sequence，
常用路线是先把 polynomial eventually bound by a smaller exponential。
例如 \(n^k\leq c\,2^n\) for sufficiently large \(n\)，然后把原式压成
\(c/2^n\)。真正写 proof 时，需要给出一个可验证的 eventually bound，
可以用 induction，也可以引用课堂已经证明过的 exponential domination。
\end{explanation}

\begin{exerciseblock}[Homework 9 Problem 2]
找出 \(L\in\R\)，并证明
\[
  \lim_{x\to0}\frac{x}{1+x}=L.
\]
\end{exerciseblock}

\begin{solution}
候选 limit 是
\[
  L=0.
\]
我们用 \(\varepsilon\)-\(\delta\) definition 证明。给定
\(\varepsilon>0\)，计算
\[
  \left|\frac{x}{1+x}-0\right|
  =
  \frac{|x|}{|1+x|}.
\]
Denominator 需要先远离 \(0\)。若 \(|x|<1/2\)，则
\[
  |1+x|\geq 1-|x|>1-\frac12=\frac12.
\]
所以在 \(|x|<1/2\) 时，
\[
  \frac{|x|}{|1+x|}
  <
  2|x|.
\]
取
\[
  \delta=\min\left\{\frac12,\frac{\varepsilon}{2}\right\}.
\]
若 \(0<|x-0|<\delta\)，则 \(|x|<1/2\)，并且
\[
  \left|\frac{x}{1+x}-0\right|
  =
  \frac{|x|}{|1+x|}
  <
  2|x|
  <
  2\delta
  \leq
  \varepsilon.
\]
因此
\[
  \lim_{x\to0}\frac{x}{1+x}=0.
\]
\end{solution}

\begin{explanation}
直接代入 \(x=0\) 会给出候选值 \(0\)，但 proof 仍然要处理 quotient。
靠近 \(0\) 时，\(1+x\) 靠近 \(1\)，所以不会接近 \(0\)。选择
\(\delta\leq1/2\) 是为了保证 \(|1+x|>1/2\)，从而把 quotient error 控制为
小于 \(2|x|\)。再选择 \(\delta\leq\varepsilon/2\)，就完成
\(\varepsilon\)-\(\delta\) estimate。
类似 \(\frac{p(x)}{q(x)}\) 的问题中，候选值来自 \(p(a)/q(a)\)，前提是
\(q(a)\neq0\)。正式 proof 的关键不是候选值本身，而是把 denominator
在 \(a\) 附近隔离开 \(0\)。通常先要求 \(|x-a|\leq\eta\)，使
\(|q(x)|\geq |q(a)|/2\)，然后把剩余 error 化成常数倍的 \(|x-a|\)。
\end{explanation}
